% From the uniformly weighted wire symbols to the tower's Laurent strip data.
\section{Weighted strip bounds for the tower kernel}
\label{sec:ex-lattice-free-kernel-estimates}

The atlas estimate in Section~\ref{sec:ex-lattice-atlas-estimates} needs a
Laurent strip bound for the tower cross-kernel and for its first residue
differences. The construction supplies these bounds from a stronger
momentum-space estimate: each entry of the alias-dependent wire symbol is
holomorphic and bounded by a summable polynomial weight in the alias. This
section proves the conversion, including the scale powers and the centered
carry in the difference. The entrywise wire estimate is then stated in its
actual form, and its source proof is cross-referenced in the audit record.

\subsection{From an analytic cube to a weighted Fourier strip}

For $p=(p_1,\ldots,p_4)\in\mathbb C^4$, let
\[
 \mathcal Q_\eta=\{p: |\operatorname{Im}p_a|<\eta,\quad
                 |\operatorname{Re}p_a|<\pi+\eta\text{ for all }a\}.
\]
Suppose $f$ is holomorphic on $\mathcal Q_\eta$, $|f(p)|\leq A$ there,
and for every real $p$ with $|p_a|<\pi+\eta$ and integer vector $n$ for
which $|p_a+2\pi n_a|<\pi+\eta$ in each coordinate,
$f(p+2\pi n)=f(p)$. Define $u$ on the horizontal strip
$|\operatorname{Im}p_a|<\eta$ by translating each real coordinate into
$[-\pi,\pi]$ before evaluating $f$. The overlap condition makes this
definition independent of the representative on the cell boundaries.
Moreover $u$ is holomorphic: on each open cell it is a translate of $f$;
on a shared boundary the two holomorphic expressions agree on a real open
box, hence agree on the complex overlap by applying the one-variable identity
theorem successively in the four coordinates. The same argument at a meeting
of several cells gives a single analytic continuation. The bound
$|u|\leq A$ follows from the bound on $f$.

Let
\[
 \widehat u(n)=(2\pi)^{-4}\int_{[-\pi,\pi]^4}
       u(x)e^{-i n\cdot x}\,dx,\qquad n\in\mathbb Z^4.
\]
For $h,b>0$ with $h+b<\eta$, shifting the $a$th contour by
$-i(h+b)\operatorname{sgn}(n_a)$ gives
\begin{equation}
 |\widehat u(n)|\leq A e^{-(h+b)|n|_1}.
 \label{eq:ex-lattice-wire-coefficient}
\end{equation}
To justify the shift, apply Cauchy's theorem on the rectangle with horizontal
sides at imaginary parts $0$ and $-(h+b)\operatorname{sgn}(n_a)$. The two
vertical integrals cancel because both $u$ and $e^{-in_ap_a}$ are
$2\pi$-periodic in the real direction. Repeat in each coordinate; the
integrand remains holomorphic and bounded throughout the product of the four
rectangles. On the shifted contour the exponential factor has modulus
$e^{-(h+b)|n|_1}$, proving the estimate.

The geometric series is bounded explicitly. Since
$e^{-b}\leq(1+b)^{-1}$,
\[
 \sum_{m\in\mathbb Z}e^{-b|m|}
 =\frac{1+e^{-b}}{1-e^{-b}}
 \leq\frac{2(1+b)}{b}.
\]
Taking products over the four coordinates and using
\eqref{eq:ex-lattice-wire-coefficient} yields
\begin{equation}
 \sum_{n\in\mathbb Z^4}e^{h|n|_1}|\widehat u(n)|
 \leq A\left(\frac{2(1+b)}{b}\right)^4.
 \label{eq:ex-lattice-wire-width}
\end{equation}
If real coefficients are needed, replace $\widehat u(n)$ by its real part;
the estimate is unchanged. If in addition
$u(-x)=\overline{u(x)}$ for real $x$, the Fourier coefficients are real,
since conjugating the integral and changing $x$ to $-x$ gives
$\overline{\widehat u(n)}=\widehat u(n)$.

\begin{lemma}[Uniform strip extraction]
\label{lem:ex-lattice-uniform-strip-extraction}
Let $\mathcal I$ be any index set, $w_i\geq0$, and let $f_i$ be a family of
symbols satisfying the preceding hypotheses with common $\eta>0$ and
$|f_i|\leq A w_i$. If $0<h,b$ and $h+b<\eta$, then their real Fourier
coefficient kernels $K_i(n)=\operatorname{Re}\widehat u_i(n)$ satisfy
\[
 \sum_{n\in\mathbb Z^4}e^{h|n|_1}|K_i(n)|
 \leq A\left(\frac{2(1+b)}b\right)^4w_i
 \quad(i\in\mathcal I).
\]
The width and the prefactor are independent of the index.
\end{lemma}

\begin{proof}
Apply \eqref{eq:ex-lattice-wire-width} to $f_i/w_i$ when $w_i>0$.
When $w_i=0$, the bound forces $f_i=0$, so the result is immediate.
Taking real parts can only decrease absolute values.
\end{proof}

\subsection{The finite-residue average and its scale powers}

For a real cross-kernel $K(r,i,j,z)$, the strip bound of width $\sigma$
and cost $A$ means
\[
 \sum_{z\in\mathbb Z^4}e^{\sigma|z|_1}|K(r,i,j,z)|\leq A
 \quad\text{for every fixed }r,i,j.
\]
The estimate is per fixed pair of finite labels; summing over labels later
contributes their cardinalities.

For an odd side $L$, let $\Lambda_L=(\mathbb Z/L\mathbb Z)^4$ and choose
the centered lift $\ell_L(r)$. Write the alias momentum as
\[
 q_l(p)=\frac{p+2\pi\ell_L(l)}{L},\qquad l\in\Lambda_L.
\]
For a word $u=(a_1,\ldots,a_s)$ in the four coordinate directions, put
\[
 \chi_u(q)=\prod_{v=1}^s(e^{iq_{a_v}}-1),\qquad \chi_{\varnothing}=1.
\]
The rescaled row symbol made from a thread $F(P,p,l)$ is
\begin{equation}
 \mathcal S_{e,u}(P,r,p)=L^{-e}\sum_{l\in\Lambda_L}
       e^{i q_l(p)\cdot\ell_L(r)}\chi_u(q_l(p))F(P,p,l).
 \label{eq:ex-lattice-spread-symbol}
\end{equation}
Here $e$ is a fixed nonnegative integer; the tower two-form uses $e=2$.
For each type pair, let $K_{e,u}(P,r,i,j,z)$ be the real part of the
normalized Fourier coefficient at $z$ of the corresponding entry of
$\mathcal S_{e,u}(P,r,\cdot)$.

Define
\[
 \Gamma_L(l)=\prod_{a=1}^4(1+|\ell_L(l)_a|),\qquad
 w_j(l)=\Gamma_L(l)^{-j}.
\]
For the centered representatives,
\begin{equation}
 \sum_{l\in\Lambda_L}\Gamma_L(l)^{-2}\leq 3^4=81.
 \label{eq:ex-lattice-residue-sum}
\end{equation}
Indeed, each coordinate lift is a distinct integer, so the left side is at
most $\prod_{a=1}^4\sum_{m\in\mathbb Z}(1+|m|)^{-2}$. The one-dimensional
sum is at most $1+2\int_1^\infty x^{-2}\,dx=3$.
Also, with $g_L(l)=1+\max_a|\ell_L(l)_a|$, we have
$g_L(l)\leq\Gamma_L(l)$ and $\Gamma_L(l)\geq1$.

Assume every scalar matrix entry of $F(P,p,l)$ is holomorphic on a fixed
complex cube $\mathcal Q_\eta$ and bounded there by
$A_0\Gamma_L(l)^{-j}$, uniformly in $P,L,l$. Assume also the overlap law
needed to glue the residue-symbol coefficients across reciprocal shifts.
For $p\in\mathcal Q_\eta$, the alias satisfies
\[
 |\operatorname{Im}q_l(p)_a|\leq\eta/L,\qquad
 |q_l(p)_a|\leq \frac{C_\eta(1+|\ell_L(l)_a|)}{L},
\]
where one may take $C_\eta=3\pi+2\eta$ for $L\geq1$.
The elementary inequality
$|e^{iz}-1|\leq |z|e^{|\operatorname{Im}z|}$ therefore gives
\[
 |\chi_u(q_l(p))|
 \leq C_\eta^s e^{s\eta}\,L^{-s}g_L(l)^s.
\]
Moreover $|e^{iq_l(p)\cdot\ell_L(r)}|\leq e^{2\eta}$ because each
$|\ell_L(r)_a|\leq (L-1)/2$. Hence, whenever $j\geq s+2$,
\begin{equation}
\begin{split}
 |\mathcal S_{e,u}(P,r,p)|
 &\leq e^{2\eta+s\eta}C_\eta^s A_0L^{-(e+s)}
       \sum_{l\in\Lambda_L} g_L(l)^s\Gamma_L(l)^{-j}\\
 &\leq 81e^{2\eta+s\eta}C_\eta^s A_0L^{-(e+s)}.
 \label{eq:ex-lattice-spread-supremum}
\end{split}
\end{equation}
The second line follows from $g_L(l)^s\leq\Gamma_L(l)^s$ and
$\Gamma_L(l)^{s-j}\leq\Gamma_L(l)^{-2}$.

Applying Lemma~\ref{lem:ex-lattice-uniform-strip-extraction} to the family
indexed by $(P,r)$ gives a common width $\sigma>0$ and a cost
$C_{e,s}L^{-(e+s)}$, where
\[
 C_{e,s}=81e^{2\eta+s\eta}C_\eta^sA_0
          \left(\frac{2(1+b)}b\right)^4
\]
for any fixed $b>0$ with $2b<\eta$. Thus each row of the cross-kernel
obtained by Fourier coefficient extraction has strip cost
$C_{e,s}/L^{e+s}$, uniformly in $P$, the torus residue, and the period.

\begin{lemma}[Residue differences are the true cover differences]
\label{lem:ex-lattice-residue-difference}
For a coefficient row $K$ on $\Lambda_L\times\mathcal C$, define
\[
 (\nabla_a^LK)(r,i,j,z)
 =K(r+\pi_L(e_a),i,j,z-c_a(r))-K(r,i,j,z),
\]
where $c_a(r)$ is the centered-lift carry characterized in the proof below.
For a word $u=(a_1,\ldots,a_s)$ set
$K_{e,\emptyset}=K_{e,0}$ and
$K_{e,(a,u)}=\nabla_a^L K_{e,u}$. If $K_{e,u}$ is the real Fourier
coefficient row of $\mathcal S_{e,u}$, then
$K_{e,(a,u)}=\nabla_a^L K_{e,u}$. Consequently the order-$s$ difference
row has width cost $C_{e,s}L^{-(e+s)}$.
\end{lemma}

\begin{proof}
The centered lift has a unique carry $c_a(r)\in\mathbb Z^4$ such that
\[
 \ell_L(r+\pi_L(e_a))+Lc_a(r)=\ell_L(r)+e_a.
\]
For odd $L$, $c_a(r)$ is $e_a$ precisely at the upper centered boundary in
coordinate $a$, and is zero otherwise. In the symbol
\eqref{eq:ex-lattice-spread-symbol}, multiplication by the new factor
$e^{iq_l(p)_a}-1$ gives
\[
 \mathcal S_{e,(a,u)}(P,r,p)
 =e^{ip\cdot c_a(r)}\mathcal S_{e,u}(P,r+\pi_L(e_a),p)
      -\mathcal S_{e,u}(P,r,p).
\]
To check the phase, use $Lq_l(p)=p+2\pi\ell_L(l)$ and
$e^{2\pi i\ell_L(l)\cdot c_a(r)}=1$; the change from
$\ell_L(r)$ to $\ell_L(r+\pi_L(e_a))$ contributes exactly
$e^{iq_l(p)_a}e^{-ip\cdot c_a(r)}$. This is the Fourier transform of
the displacement reindexing by $c_a(r)$ in the definition of
$\nabla_a^L$. Fourier coefficients are unique for these absolutely summable
rows, so the coefficient identity follows. Repeating this one-step argument
proves the formula for every word.
\end{proof}

\subsection{The actual wire receipt}

The preceding estimates apply to the construction's actual dual
cross-stencil $F(P,p,l)$. Its suppliers give the following exact data. Let
$j,t\geq4$, $k>0$, let
$\mathcal I$ be any index set, choose a flight $P_i$ of stages
$(j,t)$ and a site $l_i\in\Lambda_{L_i}$ for each $i\in\mathcal I$.
There are common $\eta,A>0$ such that every scalar matrix entry of
$F(P_i,p,l_i)$ is holomorphic on $\mathcal Q_\eta$ and bounded there by
$A\Gamma_{L_i}(l_i)^{-j}$, uniformly in the indices and the period. The
thread also obeys the precise reciprocal-shift law
\begin{equation}
 F(P,p+2\pi x,l)=F(P,p,l+\operatorname{wrap}_L(x)),
 \label{eq:ex-lattice-wire-overlap}
\end{equation}
whenever both momenta lie in the supplier's real enlarged cube. Here
$\operatorname{wrap}_L(x)$ is the residue characterized by
$\ell_L(l+\operatorname{wrap}_L(x))\equiv\ell_L(l)+x\pmod{L\mathbb Z^4}$.
The formal supplier proves these statements for every odd side $L\geq3$,
every flight with the two stage indices at least $4$, and every $k>0$.
Their factor-by-factor proof is recorded in
Section~\ref{sec:ex-lattice-free-kernel-components}; the present section
proves the quantitative consequences of those exact supplier estimates.

Here is the branch algebra in mathematical notation. Let
$m_P(p)=\operatorname{area}_{L}(p)^j$ be the patch mass factor (distinct
from the cumulative smooth factor of the train),
$b_v(p)$ the vector of four bulb factors for a train $v$, and $\beta_v(p)$
its bend factor. For a vector $g$ indexed by the four coordinate directions,
let $\mathcal M_r(g)$ be the diagonal matrix whose entry on an $r$-face is
the product of the $g_a$ over its directions. Let
\[
 \mathcal T_2(x)_{ff}=\sum_{a\in f}x_a,
 \qquad
 \operatorname{yen}(g)(p)=g(p)g(-p).
\]
For $v\in\mathbb C^4$, let $\mathsf W_r(v)$ denote the map of exterior
multiplication $\alpha\mapsto v\wedge\alpha$ from $r$-faces to
$(r+1)$-faces. Put
\[
 \Sigma_S^r(p)=\mathsf W_r\bigl((S(e^{ip_a}-1))_{a=1}^4\bigr),\quad
 \mathcal P_v(p)=\beta_v(p)^{-1}\mathcal M_2(b_v(p)^{-1}),
 \mathcal V_v(p)=\beta_v(p)\mathcal M_2(b_v(p)).
\]
Write $q_l=(p+2\pi\ell_L(l))/L$, including $q_0$ at the zero residue.
Write $r_L(q)_a=\operatorname{rail}_L(q_a)$. The $r$-flux and latch
matrices are
\[
 \Phi_{P,r}(p,l)=
  \bigl[m_P(-q_l)\mathcal M_r(r_L(-q_l))\bigr]^{\mathsf T},
 \qquad
 \Lambda_{P,r}(p)=m_P(q_0)^{-1}
       \mathcal M_r(r_L(q_0)^{-1}).
\]
The nonzero-residue port vector and scalar are
\[
 D_{P,l,a}(p)=\frac{b_{P_{\rm wing}}(q_l)_a}
                        {\operatorname{blow}_{P_{\rm wing}.mass,L}(p,l)},\qquad
R_{P,l}(p)=\operatorname{blow}_{P_{\rm wing}.mass,L}(p,l)^2
                       \beta_{P_{\rm wing}}(q_l).
\]
The land divisor is $\rho_P(p)=\operatorname{divisor}(P_{\rm land}.mass)(p)$, and the tail
vector is
\[
 \operatorname{tail}_{P,a}(p)=k^{-1}+
 \sum_{l\in\Lambda_L,\,l\ne0}
 \operatorname{yen}\!\left(m_P(\cdot)r_L(\cdot)_a\right)(q_l)
                         D_{P,l,a}(p).
\]
Further set
\[
 h_P(p)=\Lambda_{P,2}(p)\Sigma_1^1(p),\quad
 f_P(p)=\frac{\operatorname{yen}(m_P)(q_0)\,\beta_{P_{\rm land}}(p)}
                              {\beta_{P_{\rm wing}}(q_0)},\quad
 M_P(p)=\mathcal M_1(b_{P_{\rm land}}(p)),
\]
and, with $x_a(p)=\operatorname{tail}_{P,a}(p)/b_{P_{\rm land}}(p)_a$,
\[
 B_P(p)=\mathcal P_{P_{\rm wing}}(q_0)
      \bigl(\mathcal T_2(x(p))-\rho_P(p)\mathcal M_2(x(p))\bigr),
\]
\[
 N_P(p)=\rho_P(p)^2I-
 \rho_P(p)\mathcal P_{P_{\rm land}}(p)
          \Sigma_1^1(p)M_P(p)\Sigma_1^1(-p)^{\mathsf T}.
\]
For a nonzero alias $l$, let $D_{P,l}$ be the vector with components
$D_{P,l,a}$ above and set
$Y_{P,l}(p)=\operatorname{yen}(m_P)(q_l)\beta_{P_{\rm land}}(p)/R_{P,l}(p)$.
The dual stencil is exactly
\begin{align}
 W_{P,0}(p)
 &= f_P(p)\Phi_{P,2}(p,0)
   +\bigl(B_P(p)h_P(p)M_P(p)
       -\mathcal P_{P_{\rm wing}}(q_0)h_P(p)\mathcal M_1(\operatorname{tail}_P(p))\bigr)
          \Sigma_1^1(-p)^{\mathsf T}, \label{eq:ex-lattice-wire-zero}\\
 W_{P,l}(p)
 &=\mathcal P_{P_{\rm wing}}(q_l)\Sigma_L^1(q_l)
       \mathcal M_1(D_{P,l}(p))\Phi_{P,1}(p,l)\Sigma_1^1(-p)^{\mathsf T}
     +Y_{P,l}(p)\Phi_{P,2}(p,l)N_P(p),\quad l\ne0.
 \label{eq:ex-lattice-wire-nonzero}
\end{align}
These identities follow by expanding the definitions of the two branches and
the listed auxiliary matrices; they are not additional hypotheses.

The response-space notation in those formulas is as follows. For a train
$v$ and side $S$, set
$\operatorname{fieldD}_{S,v}(p)=\mathcal M_1(S^{-2}\mathscr D_v(p))$,
where $\mathscr D_v$ is the vector of level responses. The hinge matrix is
$\operatorname{hinge}_{P,r}(q)=m_P(q)\mathcal M_r(r_L(q))$; the flux is
$\Phi_{P,r}(p,l)=\operatorname{hinge}_{P,r}(-q_l)^{\mathsf T}$; and
$\operatorname{mirage}_{1,1}(p)=\Sigma_1^1(-p)^{\mathsf T}$.

Here is the exact reduction behind the branch calculations. For a train $v$
of mass side $M$, direction $a$, and momentum $p$, the level response is
\begin{align*}
 \mathscr D_{v,a}(p)&=\operatorname{feed}_{v,a}
       \bigl((\operatorname{divisor}_1)^{-1}\bigr)(p)
       +k^{-1}\operatorname{noise}_{v,a}(p),\\
 \operatorname{feed}_{v,a}(f)(p)&=M^{-2}\sum_{r\in\Lambda_M}
    D_{v,a}(q_r^{(M)}(p))D_{v,a}(-q_r^{(M)}(p))f(q_r^{(M)}(p)).
\end{align*}
The feed is the normalized finite alias sum of the two-sided drive
$D_{v,a}(q)D_{v,a}(-q)$ times its input, and the noise is the sum of the
constant-input feeds along the finite nozzle history. Every train drive is
reciprocal-periodic. The reciprocal divisor input is periodic because its
chip factors are invariant under integer reciprocal shifts. A finite alias
sum of a periodic summand is periodic: a reciprocal shift permutes its alias
sites. More explicitly, for a hop of side $L$ from a train of mass side $M$,
the mixed-radix alias bijection $\Lambda_{LM}\to\Lambda_L\times\Lambda_M$
is determined by
$\ell_{LM}(r)\equiv\ell_L(l)+L\ell_M(v)\pmod{LM}$, with the unique
centered representatives differing by a multiple of $LM$ absorbed into a
reciprocal shift. This is a bijection because every residue modulo $LM$ has
a unique decomposition into a residue modulo $L$ and then a residue modulo
$M$ after division by $L$. Thus
$q_r^{(LM)}(p)=q_v^{(M)}(q_l^{(L)}(p))$ modulo $2\pi\mathbb Z^4$.
Reindexing this bijection in the normalized feed sum and using
$(LM)^{-2}=L^{-2}M^{-2}$ gives
\[
 \operatorname{feed}_{\mathrm{hop}}(f)
   =\operatorname{slide}_{L}(\operatorname{feed}_{v}(f)),\qquad
 D_{\mathrm{hop},a}(z)=\operatorname{pullLine}_{P,a}(Mz)D_{v,a}(z).
\]
The nozzle list has a top entry with constant drive, contributing $1$;
each earlier entry is the lifted nozzle list of the predecessor train.
Reindexing the same aliases gives
$\operatorname{Noise}_{\mathrm{hop}}=1+\operatorname{slide}_{L}
\operatorname{Noise}_{v}$. Consequently
$\mathscr D_{\mathrm{hop},a}=\operatorname{slide}_{L,a}
\mathscr D_{v,a}+k^{-1}$, the response feedback used below. Induction on
the finite nozzle length proves reciprocal periodicity of every
$\mathscr D_{v,a}$.
The level scalar
\[
 \rho_v(p)=\sum_a\frac{J_a(p)}{\mathscr D_{v,a}(p)},\qquad
 J_a(p)=(e^{ip_a}-1)(e^{-ip_a}-1),
\]
is periodic because each $J_a$ and each $\mathscr D_{v,a}$ is periodic.
Off $p=0$, the source recovery identity is
$\rho_v=\operatorname{divisor}_M(p)^2\beta_v(p)$ and the level response
identity is $\mathscr D_{v,a}=b_{v,a}/\operatorname{divisor}_M(p)$. Substituting
these identities in the definitions of pulse and vault gives
\[
 \operatorname{pulse}_v(p)=\rho_v(p)^{-1}
        \mathcal M_2(\mathscr D_v(p)^{-1}),\qquad
 \operatorname{vault}_v(p)=\rho_v(p)\mathcal M_2(\mathscr D_v(p)).
\]
Thus pulses are reciprocal-periodic on every nonzero overlap. A nonzero
reciprocal shift cannot have both endpoints in the small real cube while one
endpoint is zero; the elementary overlap separation inequality proves this.
These are the level-response and pulse identities used below, not
periodicity of the unscaled divisor $\operatorname{divisor}_M$ (which is
not periodic when $M>1$).

For clarity, the branch algebra is verified here. All products below are
matrix products in the indicated face spaces. At a nonzero alias write
$q=q_l$, $s=\rho_P(p)$, $D=\mathcal M_1(D_{P,l})$, $D'=
\mathcal M_1(b_{P_{\rm land}}(p))$, and abbreviate the two order-one
signals and the order-one/order-two fluxes by their displayed symbols. The
port definitions and the divisor recovery identity give
\[
\operatorname{chainY}_P(p,q)
   =Y_{P,l}(p)\,\rho_P(p)^2,
\qquad
\operatorname{portD}_{P,l,a}(p)=L^{-2}\mathscr D_{P_{\rm wing},a}(q_l).
\]
The first equality follows by substituting
$R_{P,l}=\operatorname{blow}_{P_{\rm wing}.mass,L}^2
\beta_{P_{\rm wing}}(q)$ into $Y_{P,l}$ and then using
$\rho_P^2\beta_{P_{\rm land}}/\beta_{P_{\rm wing}}$ for the ratio
of the two bend factors; the second is the definition of the port vector
after cancellation of its blow-up denominator. Also
\[
 N_P=s^2 I-s\,\mathcal P_{P_{\rm land}}(p)\Sigma_1^1(p)D'
                   \Sigma_1^1(-p)^{\mathsf T},
 \qquad
 \operatorname{notch}_P=s^2(I-\operatorname{cloud}_P).
\]
For the first identity use
$s\operatorname{fieldD}_{1,P_{\rm land}}=\mathcal M_1(b_{P_{\rm land}})$
in the definition of the cloud; the second follows from the mouth identity
and the definition of the notch. Substituting these identities in the
nonzero branch, and using $Y_{P,l}s^2=\operatorname{chainY}_P$, gives
exactly
\[
 \mathcal P_{P_{\rm wing}}(q)\Sigma_L^1(q)D\Phi_{P,1}
       \Sigma_1^1(-p)^{\mathsf T}
 +\operatorname{chainY}_P(p,q)\Phi_{P,2}(p,l)
       (I-\operatorname{cloud}_P(p))
 =W_{P,l}(p).
\]
This is the stake branch identity, including its order of matrix factors.

The heel is, by definition,
\[
 \operatorname{heel}_{P,a}(p)=
   \operatorname{yen}(\operatorname{pullLine}_{P,a})(q_0)
       b_{P_{\rm wing},a}(q_0).
\]
The relevant one-step identity for this response is
\[
 \mathscr D_{P_{\rm land},a}(p)=\operatorname{slide}_{P,a}
      \mathscr D_{P_{\rm wing},a}(p)+k^{-1},
 \qquad
 (\operatorname{slide}_{P,a}g)(p)=L^{-2}\sum_{l\in\Lambda_L}
      \operatorname{yen}(\operatorname{pullLine}_{P,a})(q_l)g(q_l).
\]
Use the port-to-response equality
$\operatorname{portD}_{P,l,a}=L^{-2}\mathscr D_{P_{\rm wing},a}(q_l)$
in this finite sum. The zero-residue summand is
$\operatorname{heel}_{P,a}(p)/s$, by the zero-port identity
$\operatorname{blow}_{P_{\rm wing}.mass,L}(p,0)=s$; the remaining
residues sum to the above-alias tail. Adding $k^{-1}$ gives
$\mathscr D_{P_{\rm land},a}=\operatorname{heel}_{P,a}/s+
\operatorname{tail}_{P,a}$. Since off the origin
$\mathscr D_{P_{\rm land},a}=b_{P_{\rm land},a}/s$, multiplication by $s$
proves the bulb network identity
\[
 b_{P_{\rm land},a}(p)=\operatorname{heel}_{P,a}(p)+
             s\operatorname{tail}_{P,a}(p).
\]
For the zero branch write $q=q_0$, $s=\rho_P(p)$, $H$ and $H'$
for the wing and land pulses at $q$ and $p$, $d=\Sigma_1^1(p)$,
$\delta=\Sigma_L^1(q)$, $v=\Phi_{P,2}(p,0)$,
$u=\Phi_{P,1}(p,0)$, and $m=\operatorname{mirage}_{1,1}(p)$.
Let $f_1=\operatorname{hinge}_{P,1}(q)$ and
$f_2=\operatorname{hinge}_{P,2}(q)$, $\lambda=\Lambda_{P,2}(p)$,
$R=\mathcal M_1(\operatorname{tail}_P(p))$, $B=B_P(p)$,
$D=\operatorname{fieldD}_{L,P_{\rm wing}}(q)$,
$D'=\operatorname{fieldD}_{1,P_{\rm land}}(p)$, and
$y=\operatorname{foot}_P(p)$. These symbols will be used in the brace,
one-form, and cancellation identities below.
For the brace relation, abbreviate the land and wing bulb vectors at their
respective arguments by $f$ and $g$, the two patch masses by
$c_0=m_P(q_0)$ and $d_0=m_P(-q_0)$, and the corresponding rails by
$E_a=\operatorname{rail}_L(q_{0,a})$ and
$F_a=\operatorname{rail}_L(-q_{0,a})$. Put
$Z_a=\operatorname{heel}_{P,a}(p)/f_a$ and
$x_a=\operatorname{tail}_{P,a}(p)/f_a$. The definitions of heel and the
port-to-response formula give, coordinate by coordinate,
\[
 g_a^{-1}Z_a=c_0d_0\,F_a f_a^{-1}E_a.
\]
Inserting this in the diagonal two-form pulse and multiplying by the land
pulse, two-form flux, and two-form hinge gives
$y\,vH'f_2=H\mathcal M_2(Z)$. The bulb network identity gives
$Z_a=1-sx_a$. On a two-face $\{a,b\}$, direct multiplication yields
\[
 1-(1-sx_a)(1-sx_b)=s(x_a+x_b-sx_ax_b).
\]
Since $\mathcal T_2(x)$ has diagonal entry $x_a+x_b$ and
$\mathcal M_2(x)$ has entry $x_ax_b$, this is the matrix identity
$I-\mathcal M_2(1-sx)=s(\mathcal T_2(x)-s\mathcal M_2(x))$.
Multiplying by $H=\mathcal P_{P_{\rm wing}}(q_0)$ proves
$sB=H-yvH'f_2$, the brace relation used in the cancellation.

For the one-form flux, each matrix entry of the hinge-flux product is
diagonal in the coordinate direction. At the zero alias it equals
\[
 \bigl(f_1D u\bigr)_{aa}
 =\operatorname{yen}(m_P)(q_0)\operatorname{rail}_L(q_{0,a})
   \operatorname{rail}_L(-q_{0,a})
   \operatorname{fieldD}_{L,P_{\rm wing}}(q_0)_a.
\]
Because $\operatorname{fieldD}_{L,P_{\rm wing}}(q_0)_a=
b_{P_{\rm wing},a}(q_0)/s$,
the right side is $\operatorname{heel}_{P,a}(p)/s$. The scalar toe-feed
identity above says this equals
$\mathscr D_{P_{\rm land},a}(p)-\operatorname{tail}_{P,a}(p)$, proving
$f_1D u=D'-R$ entry by entry. The equality
$\operatorname{chainY}_P(p,q_0)=\operatorname{foot}_P(p)$ follows as
follows. Divisor recovery gives
$\rho_{P_{\rm land}}=s^2\beta_{P_{\rm land}}$ and
$\rho_{P_{\rm wing}}=d_w^2\beta_{P_{\rm wing}}$, where
$d_w=\operatorname{divisor}(P_{\rm wing}.mass)(q_0)$. The zero-port
identity gives $s=L^2d_w$. Since
$\operatorname{scaleR}_{1,P_{\rm land}}=\rho_{P_{\rm land}}$ and
$\operatorname{scaleR}_{L,P_{\rm wing}}=L^4\rho_{P_{\rm wing}}$,
their ratio is $\beta_{P_{\rm land}}/\beta_{P_{\rm wing}}$. Multiplying by
$\operatorname{yen}(m_P)(q_0)$ gives both the chain and foot expressions.

The bank and land-response identities are
\[
 \operatorname{mouth}_P=\mathcal M_1(b_{P_{\rm land}}),\qquad
 s\operatorname{fieldD}_{1,P_{\rm land}}=\mathcal M_1(b_{P_{\rm land}}).
\]
Thus $\operatorname{mouth}_P=sD'$.
The latch inverse, hinge transport, and the order-one flux calculation give
the identities
\begin{align*}
 f_2\lambda=I,\qquad \lambda f_2=I,\\
 f_2\delta&=d f_1,\\
 f_1D u&=D'-R,\\
sB&=H-y\,vH'f_2,\qquad \operatorname{mouth}_P=sD'.
\end{align*}
The hinge-transport line follows entrywise from the alias relation
$L\operatorname{rail}_L(q_{0,a})\operatorname{chip}(q_{0,a})
=\operatorname{chip}(p_a)$, inserted into the definitions of the hinge and
signal. The third line is the one-form identity obtained by
multiplying the flux by its hinge; the fourth is the brace identity after
the two-form flux calculation. The mouth identity was derived just above.
Also $\operatorname{chainY}_P(p,q)=y$. These identities imply, with no
commutation of noncommuting factors,
\begin{align*}
 H\delta D u\,m
   &=H\lambda d f_1D u\,m
     =H\lambda d(D'-R)m,\\
 B\lambda d\operatorname{mouth}_P\,m-H\lambda dR\,m
   &=H\lambda dD'm-yvH'dD'm-H\lambda dR\,m.
\end{align*}
For the first equality use $\delta=\lambda d f_1$ (which follows from
$f_2\delta=df_1$ and $f_2\lambda=I$). For the second use
$\operatorname{mouth}_P=sD'$ and $sB=H-yvH'f_2$, then cancel the adjacent
$f_2\lambda$. Adding $yv$ to the second line yields
$H\lambda d(D'-R)m+yv(I-H'dD'm)$, and its first term is
$H\delta D u\,m$. This is precisely the zero-alias script formula and
proves the copper branch identity. Each displayed identity is an equality
of the finite matrices defined above; the component chapter supplies the
scalar divisor, rail, and bend inputs on the common analytic cube.

The reciprocal-shift law used in the Fourier argument also follows from
the actual wire symbol, rather than being an independent periodicity
assumption. For $x\in\mathbb Z^4$ and a residue $l$, the centered alias
arithmetic gives a cover vector $y$ such that
\[
 q_l(p+2\pi x)=q_{l+\operatorname{wrap}_L(x)}(p)+2\pi y.
\]
Every factor of the script is invariant under these reciprocal shifts.
The level responses are periodic by the finite-feed and nozzle induction
above; hence their diagonal matrices and the scalar $\rho_v$ are periodic.
The pulses are periodic on nonzero overlaps by their displayed level-source
formula. The signals and the jewel factors are periodic because they are
formed from $e^{ip_a}-1$ and its opposite-momentum partner. The patch mass is
periodic by the plank identity. The scalar factor is exactly
\[
 \operatorname{chainY}_P(p,q)=
 \operatorname{yen}(m_P)(q)\,
 \frac{\rho_{P_{\rm land}}(p)}{L^4\rho_{P_{\rm wing}}(q)}.
\]
It is periodic by the patch-mass and level-response identities. The other
factors in $\operatorname{scriptT}$ are level responses, pulses, signals,
hinges, or fluxes already shown periodic. Substitution in its definition
therefore yields
$\operatorname{scriptT}_P(p+2\pi x,q+2\pi y)=
\operatorname{scriptT}_P(p,q)$. For nonzero external momentum, the copper and
stake branch identities just proved identify the wire with this script;
the zero momentum endpoint follows directly from the definitions. Combining
this script invariance with the residue permutation
$l\mapsto l+\operatorname{wrap}_L(x)$ gives
\eqref{eq:ex-lattice-wire-overlap}. More explicitly, when both external
momenta lie in the real enlarged cube and $x\ne0$, both are nonzero; the
alias exchange above supplies the compensating fine reciprocal vector, and
the shift identity for the script removes it. When $x=0$ the claim is
identity. This is the reciprocal-shift content of the wire matching
argument; the companion audit records its exact source declarations.
The real enlarged cube has coordinate radius $\pi+\mathrm{zone}$, where
$\mathrm{zone}=\pi/4$ is the padding. Since $\pi+\mathrm{zone}<2\pi$,
no nonzero reciprocal vector can connect zero momentum to another point
of this cube; this is the separation used above.

The product estimates use the following weights, all uniform on a common
complex cube and independent of $L$:
\begin{center}
\begin{tabular}{c|c}
factor & weight in the alias label $l$ \\
\hline
mass factor $m_P(q_l)$ & $\Gamma_L(l)^{-j}$ \\
rail factors and reciprocal rails at the zero alias & $1$ \\
dual multipliers, bend ratios at zero alias, correction and tail terms & $1$ \\
 scaled signal $\Sigma_L^1(q_l)$ & $g_L(l)$ \\
 port mask $\mathcal M_1(D_{P,l})$ & $g_L(l)^{-2}$ \\
flux $\Phi_{P,r}(p,l)$ & $\Gamma_L(l)^{-j}$ \\
levy and notch at nonzero alias & $1$
\end{tabular}
\end{center}
The component suppliers prove these bounds by the estimates listed in
Section~\ref{sec:ex-lattice-free-kernel-components}. For the
zero branch every product has total weight $1$. In the first nonzero product
the total weight is at most
$g_L(l)g_L(l)^{-2}\Gamma_L(l)^{-j}
 \leq\Gamma_L(l)^{-j}$; the second has weight
$\Gamma_L(l)^{-j}$. For matrix products, apply
$|(AB)_{ij}|\leq\sum_r|A_{ir}||B_{rj}|$; every finite intermediate type
sum contributes a fixed cardinality. Taking the larger of the two branch
constants and the smaller of their widths yields the entrywise weighted
analytic estimate stated above. The branch multiplication and the weights it
must deliver are explicit here; Section~\ref{sec:ex-lattice-free-kernel-components}
derives the component estimates used in that multiplication.

\begin{theorem}[Wire symbol yields the tower's order-$e$ strip receipt]
\label{thm:ex-lattice-wire-tower-strip}
Under $j\geq4$, odd $L$, and the entrywise weighted analytic bound together
with the reciprocal-shift overlap identity \eqref{eq:ex-lattice-wire-overlap},
there are $\sigma>0$ and $C_{e,j}>0$, independent of the flight,
residue, and period, such that for every word $u$ of length $s$ with
$s\leq j-2$, the cross-kernel row $K_{e,u}$ has strip width $\sigma$ and cost
$C_{e,j}/L^{e+s}$ for every flight, residue, and fixed type pair.
In particular $e=2$ gives the tower component $T$ the base cost
$C_{2,j}/L^2$ and each first-difference cost $C_{2,j}/L^3$ required by
Lemma~\ref{lem:ex-lattice-cross-atlas}.
\end{theorem}

\begin{proof}
For a fixed word $u$ of length $s$, the finite-residue estimate
\eqref{eq:ex-lattice-spread-supremum} gives a common analytic bound with
factor $L^{-(e+s)}$. To check the overlap, let $x\in\mathbb Z^4$ and let
$l'=l+\operatorname{wrap}_L(x)$. By the defining congruence for the wrap,
there is $n\in\mathbb Z^4$ such that
\[
 q_l(p+2\pi x)=q_{l'}(p)+2\pi n.
\]
The chip factors and the exponential displacement factor are unchanged by
this integral reciprocal shift, and the supplier law
\eqref{eq:ex-lattice-wire-overlap} changes $F(P,p+2\pi x,l)$ into
$F(P,p,l')$. The map $l\mapsto l'$ is a permutation of the finite residue
set, so reindexing the finite sum proves equality of the two row symbols on
the real overlap. Holomorphic uniqueness extends the equality to the
overlapping complex domain used for coefficient extraction.
Lemma~\ref{lem:ex-lattice-uniform-strip-extraction}
therefore gives the strip cost $C_{e,s}/L^{e+s}$. There are only finitely
many words in the four coordinate directions of length at most $j-2$.
The width $\sigma$ is the same for every word; take the maximum of the
finitely many costs $C_{e,s}$ over all words of length at most $j-2$ to
obtain one cost $C_{e,j}$. Finally
Lemma~\ref{lem:ex-lattice-residue-difference}, iterated along $u$, identifies
the word's Fourier coefficient row with the iterated cover difference.
This is exactly the stated receipt.
\end{proof}

\begin{remark}[Remaining local analytic inputs]
This proof discharges the passage from a uniform alias-weighted analytic
bound to all required tower strip receipts, with explicit $L^{-2-s}$
scaling and the carry-corrected differences. Its mathematical dependence on
the actual dual-kernel producer is restricted to the entrywise weighted
analytic estimate and the component inequalities listed in the audit. In particular, the
finite-residue geometry, scale powers, geometric coefficient summation,
overlap gluing, and finite uniformization are proved here. The selected
response atlas then follows from Theorem~\ref{thm:ex-lattice-response-atlas}.
\label{rem:ex-lattice-free-kernel-inputs}
\end{remark}
